\SourceLeaf{376}%
Zwecke. Um 1811 war ein Dreiecksnetz von Fiume bis Turin fertiggestellt. Es erstreckte sich somit in Richtung eines Parallelkreises auf dem 45. Breitengrad über sieben Längengrade. Um diese Arbeiten auch für die weitergehende wissenschaftliche Erforschung der Erdgestalt zu nutzen, beschloss die damalige französische Regierung auf \textsc{Laplace}’ Vorschlag, das Dreiecksnetz nach Westen bis zum Atlantik zu erweitern und die für eine Längengradmessung erforderlichen Arbeiten damit zu verbinden. Die Arbeit wurde sofort mit Eifer begonnen, durch die damaligen Ereignisse für einige Zeit unterbrochen und bald mit gleichem Einsatz fortgesetzt. 1818 war das Dreiecksnetz über das französische Gebiet vom Atlantik bei Bordeaux bis an die Grenze Savoyens vermessen. Zur Vollendung des geodätischen Teils fehlte damit nur noch das Stück im Herrschaftsgebiet des Königs von Sardinien. Die dortige und die österreichische Regierung erkannten beide die wissenschaftliche Bedeutung dieses großen Unternehmens mit Nachdruck an und beschlossen, eine Kommission aus Astronomen und Offizieren beider Staaten mit den noch ausstehenden geodätischen Arbeiten und den in Italien erforderlichen astronomischen Arbeiten zu beauftragen. Diese Arbeiten bilden den Inhalt des vorliegenden Werks, das, wie es scheint, von den Astronomen \textsc{Carlini} und \textsc{Plana} gemeinsam redigiert wurde.

Der erste Band ist ausschließlich den geodätischen Arbeiten gewidmet. Die beiden östlichen Endpunkte des Dreiecksnetzes in Frankreich, der Mont Colombier und der Mont Granier (unweit von Chambery), bilden die Seite, an der die neue Messung ansetzen musste. Sie war bis zur westlichsten Seite des Netzes in der Lombardei, Massé --- Superga (bei Turin), fortzuführen. In diesem Gelände mit den höchsten Gebirgen Europas hätte man erwarten können, leicht große Dreiecke bilden zu können und mit sehr wenigen Zwischenpunkten auszukommen --- der Mont Granier liegt nur 150000 Meter von Superga entfernt. Doch gerade auf dieser mäßig langen Strecke hatte man im Gegenteil mit den größten Schwierigkeiten zu kämpfen. Die höheren Gipfel sind gar nicht oder nur schwer zugänglich, die Baustoffe für die Signale lassen sich nur mit größter Anstrengung hinaufschaffen, und heftige Stürme bedrohen
